\begin{textsample}{POS Dim 4 – human – Score 4.00 – mg/mg\_ef\_linguagens\_ingles\_9\_p1.txt}  \label{ex:f4_pos_017}
Pesquisas recentes indicam que a \textbf{leitura}, longe de ser uma atividade passiva, é um processo dinâmico através do qual o leitor se envolve ativamente na (re)criação do sentido do \textbf{texto}, fundamentado não só em seu conhecimento anterior, mas também nas condições de produção \textbf{textual} ( o que foi escrito, por/para quem, com que propósito, de que forma, quando e onde ).
A \textbf{leitura} pode ser \textbf{vista} como uma interação à distância entre leitor e autor através do \textbf{texto}, inserida nas situações comunicativas do cotidiano.
Os três tipos de conhecimento ( o de mundo, o linguístico e o \textbf{textual} ) são mobilizados nessa relação dialógica leitor-autor, envolvendo também o uso de estratégias de \textbf{leitura} num processo contínuo de predizer, antecipar, formar e reformular hipóteses, avaliar, concordar, discordar, inferir e “ ler nas entrelinhas ” com base em pistas \textbf{textuais} e contextuais.

% matched lemmas: leitura, texto, textual, ver
\end{textsample}
